\section{自己開示の受容}\label{受容}

前節で述べたように，自己開示が開示者と被開示者の間での相互作用であり，特に自己開示の相互性のプロセスにおいて被開示者の受容が重要な役割を果たすと述べた．そこで本節では，自己開示の受容に焦点を当て，自己開示の促進と，人間関係の発展に果たす役割と影響について説明する．

% --------------------------------------------------------------------------------------------------------

\subsection{受容の定義}

被開示者の反応は，受容と拒絶の2つに分けられる．森脇\cite{森脇_2005}によると，被開示者の受容には「聴く際の真摯な姿勢」「アドバイス」「親身な行動」「共感」，拒絶には「否定・無視」「無関心」「聴く際の真剣味のなさ」「無反応」といった多様な反応が含まれる．これらを踏まえ，川西は，相手の自己開示に対する受容を「開示者に対する肯定的・共感的反応」，拒絶を「開示者に対する否定的・非共感的反応」と定義している．


\subsection{自己開示の促進に与える影響}\label{}

受容は自己開示の促進に重要な役割を果たすことが示されている．相川\cite{相川充_2010}によれば，相手の共感的な反応を通じて受容されていると感じることで，開示者はより内面的な自己開示へと動機づけられるという．この過程について，Weber\cite{weber_2004}の研究では2つの重要な知見が示されている．第一に，親密な相手から思いやりや共感を受けることで，相手に理解されているという感覚が高まることと，第二に，親密な相手からの思いやりや共感を受けることは，相手に対して意図的に自己開示を行おうとする傾向と関連があることが示されている．
% ただし，実際の自己開示の量や深さとの関連は見られなかったことから，思いやりや共感の表現は，自己開示の行動そのものよりも，開示しようとする意図に影響を与える可能性が示唆されている．

対照的に，自己開示に対する拒絶的な反応は，開示者に心理的な傷つきをもたらし，その後の自己開示を抑制する可能性がある．拒絶的な反応を経験した個人は，その後の人間関係において防衛的になり，深い自己開示を避ける傾向があることが指摘されている．

このように，被開示者から受容されたという感覚は，さらなる自己開示を促進すると考えられる．また，受容的反応は自己開示の促進だけでなく，関係の親密化においても重要な役割を果たすことが明らかになっている．次項では，受容が関係の親密化に与える影響について説明する．



\subsection{関係の親密化に与える影響}\label{}

前項で述べたように，受容は自己開示を促進するだけでなく，関係の親密化においても重要な役割を果たすことが明らかになっている．榎本は，自己開示に対する受容が開示者の自尊心を高め，2者間の人間関係が親密化することを指摘している．

この関係性について，Reis\cite{reis_2017}は理解（understanding）という概念から詳細な説明を行っている．理解とは相手の経験や特性についての知識や認識を持つことであり，理解されていると感じることは関係性の質を高めることが示されている．さらに，Weberの研究によれば，思いやりや共感の表現は，お互いの信頼関係を深める要因となることが示されている．

相手からの思いやりや共感を受けることで相手に理解されているという感覚が高まり，それが結果として関係性の質の向上につながると考えられる．これらの知見は，自己開示に対する受容的な反応が，単に自己開示を促進するだけでなく，関係の親密化に重要な役割を果たしていることを示している．


\subsection{自己開示に対する受容の課題}\label{受容の課題}
% ここまでは受容の重要性を述べた → その難しさ → それを解決する必要あり(?)
これまで述べてきたように，自己開示の受容は人間関係の親密化において重要な役割を果たす．しかし，自己開示に対する適切な受容は必ずしも容易ではない．特に深い自己開示を受けた際，被開示者は適切な反応を示すことに困難を感じることがある．
表層的な自己開示に対しては比較的容易に受容できるが，深い自己開示に対しては，より慎重で共感的な反応が必要となる．

% 受容側に着目した研究少ない．やる必要あるって話
川西が，従来の研究では被開示者の受容が前提とされる場合の効果に関するものがほとんどであると指摘しているように，既存研究では，開示者側に着目した研究が多く，被開示者の受容に着目したものは比較的少ない．しかし，自己開示した際，相手に必ずしも受け入れられるとは限らない．特に，深い自己開示は過去のトラウマや失敗など，否定的側面を含むことが多い．開示内容が否定的である場合，批判や誤解を招く可能性や無関心な態度を取られる可能性が十分にある．
